% PASSAGE Paper - arXiv Preprint
% Comprehensive first draft based on PUBLICATION.md skeleton
\documentclass{article}

% arXiv preprint style (can be swapped for conference style later)
\usepackage[preprint]{arxiv}  % Options: [preprint], [final]

% Core packages
\usepackage[utf8]{inputenc}
\usepackage[T1]{fontenc}
\usepackage{hyperref}
\usepackage{url}
\usepackage{booktabs}
\usepackage{amsfonts}
\usepackage{amsmath}
\usepackage{amssymb}
\usepackage{nicefrac}
\usepackage{graphicx}
\usepackage{xcolor}
\usepackage{subcaption}
\usepackage{algorithm}
\usepackage{algorithmic}
\usepackage{multirow}
\usepackage{siunitx}
\usepackage{enumitem}

% Code listings
\usepackage{listings}
\lstset{
  basicstyle=\ttfamily\small,
  breaklines=true,
  frame=single,
  language=Python,
  keywordstyle=\color{blue},
  commentstyle=\color{gray},
  stringstyle=\color{orange},
  numbers=left,
  numberstyle=\tiny\color{gray},
}

% Custom commands
\newcommand{\dataset}{\textsc{Passage}}
\newcommand{\ie}{\textit{i.e.}}
\newcommand{\eg}{\textit{e.g.}}
\newcommand{\cf}{\textit{cf.}}
\newcommand{\etal}{\textit{et al.}}

% Title
\title{PASSAGE: A Multi-Resolution Elevation-Based Pathfinding Dataset\\for Certified Machine Learning Systems}

% Authors
\author{%
  Author One\thanks{Corresponding author. Equal contribution.} \\
  Affiliation One\\
  \texttt{author1@example.com} \\
  \And
  Author Two\footnotemark[1] \\
  Affiliation Two\\
  \texttt{author2@example.com} \\
}

\begin{document}

\maketitle

% ============================================================================
% ABSTRACT
% ============================================================================
\begin{abstract}
We present \dataset{}, a large-scale, multi-resolution pathfinding benchmark derived from the JAXA ALOS AW3D30 Digital Surface Model. Unlike existing grid-based benchmarks that rely on binary obstacles and fixed resolutions, \dataset{} provides continuous elevation-derived cost maps, procedurally generated superellipse obstacles with configurable shape distributions, and ground-truth paths computed by multiple solver backends. The dataset spans seven resolution levels (64$\times$64 to 4096$\times$4096 pixels) with 13,875 samples in repository defaults, organized into stratified train/calibration/validation/test splits. We provide five physically-motivated cost models (elevation, distance, slope-uphill, energy, absolute-elevation) and three solver backends (NetworkX, pathfinding, NumPy-grid) supporting A*, Dijkstra, and bidirectional variants. Designed with aerospace certification standards (ARP6983, ED-324) in mind, \dataset{} offers deterministic generation, complete provenance metadata, and a dedicated calibration split for certification-oriented experiments. The toolkit includes direct Hugging Face Hub export with zstd-compressed shards and configurable chunk sizes. We release the full generation pipeline, configuration utilities, and baseline experiments to enable reproducible research in supervised surrogate modeling, reinforcement learning, and graph neural networks for pathfinding. The dataset and code are available at \url{https://huggingface.co/datasets/passage}.
\end{abstract}

% ============================================================================
% 1. INTRODUCTION
% ============================================================================
\section{Introduction}
\label{sec:introduction}

\subsection{Context and Motivation}

The proliferation of learning-based approaches for navigation and pathfinding---including supervised surrogate models, reinforcement learning (RL) agents, and graph neural networks (GNNs)---has created substantial demand for standardized benchmark datasets. Such benchmarks are essential for fair comparison of algorithms, reproducible research, and systematic evaluation of generalization capabilities.

However, existing pathfinding benchmarks suffer from several critical limitations. Grid-based benchmarks from the Moving AI Lab~\cite{sturtevant2012benchmarks} provide maps derived from video games (Dragon Age, StarCraft) but feature only binary obstacles and fixed resolutions, lacking the continuous cost gradients found in real terrain. Random grid environments offer scalability but sacrifice terrain realism. Continuous navigation datasets from robotics and autonomous driving (\eg, nuScenes, Waymo) incorporate sensor modalities and dynamics but are unsuitable for evaluating discrete grid-based algorithms.

Simultaneously, safety-critical domains such as aerospace, autonomous vehicles, and robotics increasingly require \emph{certifiable} AI systems operating within well-defined Operational Design Domains (ODD). Emerging standards like SAE ARP6983~\cite{arp6983} and EUROCAE ED-324~\cite{ed324} mandate reproducible test conditions, bounded performance metrics, and evidence of data quality---requirements that existing pathfinding datasets do not adequately address.

\subsection{Problem Statement}

We identify the following gaps in current pathfinding datasets:
\begin{itemize}[noitemsep]
    \item \textbf{Limited resolution variety}: Most benchmarks offer single or few fixed resolutions, preventing studies of computational scalability and resolution-dependent generalization.
    \item \textbf{Binary obstacles only}: Existing datasets lack continuous terrain-based cost models that capture elevation, slope, and energy expenditure.
    \item \textbf{Insufficient sample counts}: Many benchmarks provide too few samples for modern deep learning architectures requiring thousands of training examples.
    \item \textbf{No certification alignment}: Current datasets lack the provenance metadata, deterministic generation, and split structures needed for certification workflows.
    \item \textbf{Single solver ground truth}: Most datasets provide paths from a single algorithm, preventing studies of solver agreement and algorithm selection.
\end{itemize}

\subsection{Contributions}

We address these gaps with the following contributions:

\begin{enumerate}[noitemsep]
    \item \textbf{\dataset{} Dataset}: A large-scale, multi-resolution pathfinding benchmark with 13,875 samples (repository defaults) across seven resolutions (64$\times$64 to 4096$\times$4096), derived from real-world elevation data (JAXA ALOS AW3D30).

    \item \textbf{Procedural Obstacle Generation}: Configurable superellipse-based obstacles (Lamé curves) with flexible shape distributions (log-uniform, uniform, beta, exponential), size ratios, rotation angles, and coverage constraints up to 30\%.

    \item \textbf{Multi-Backend Solver Framework}: Three solver backends (NetworkX, pathfinding, NumPy-grid) implementing A*, Dijkstra, bidirectional variants, and Bellman-Ford, with five physically-motivated cost models.

    \item \textbf{Certification-Oriented Design}: Bounded ODD definition, deterministic generation with configurable seeds, complete provenance metadata per sample, and a dedicated calibration split (10\%) for conformal prediction and uncertainty quantification experiments.

    \item \textbf{Open-Source Toolkit}: Complete pipeline from tile download through Hugging Face Hub export, with zstd-compressed shards and configurable per-resolution chunk sizes for efficient storage and streaming.
\end{enumerate}

\subsection{Paper Organization}

Section~\ref{sec:related} reviews related work on pathfinding datasets, learning-based approaches, and certification standards. Section~\ref{sec:dataset} presents the dataset design, including data sources, obstacle generation, and cost models. Section~\ref{sec:pipeline} describes the generation pipeline and export mechanisms. Section~\ref{sec:solvers} details solver implementations and evaluation metrics. Section~\ref{sec:experiments} presents experimental validation including solver benchmarks and baseline learning models. Section~\ref{sec:research} proposes research directions enabled by \dataset{}. Section~\ref{sec:conclusion} concludes with future work.

% ============================================================================
% 2. RELATED WORK
% ============================================================================
\section{Related Work}
\label{sec:related}

\subsection{Existing Pathfinding Datasets}

\paragraph{Grid-Based Benchmarks.}
The Moving AI Lab benchmarks~\cite{sturtevant2012benchmarks} remain the most widely used pathfinding testbeds, providing maps from Dragon Age, StarCraft, and other games. While valuable for algorithm comparison, these maps feature binary obstacles (passable/impassable) and fixed resolutions, limiting their applicability to terrain-aware planning. Random grid generators offer scalability but lack geographic realism.

\paragraph{Continuous Environment Datasets.}
Robot navigation datasets~\cite{mirowski2017learning} and autonomous driving benchmarks (nuScenes~\cite{caesar2020nuscenes}, Waymo~\cite{sun2020waymo}) operate in continuous state spaces with rich sensor modalities. However, their focus on perception and dynamics makes them unsuitable for evaluating discrete grid-based pathfinding algorithms.

\paragraph{Elevation-Based Data.}
Digital Elevation Models (DEMs) from sources like SRTM, ASTER, and ALOS are available for GIS applications but lack standardized pathfinding ground truth, train/test splits, and obstacle annotations suitable for ML evaluation.

\subsection{Learning-Based Pathfinding}

\paragraph{Supervised Surrogate Modeling.}
Convolutional neural networks have been applied to path prediction, with U-Net architectures~\cite{ronneberger2015unet} producing path segmentation masks. Encoder-decoder models estimate cost-to-go functions, while attention mechanisms capture long-range dependencies. However, these approaches typically assume fixed input resolutions and struggle with generalization across terrain types.

\paragraph{Value Iteration Networks.}
Tamar \etal~\cite{tamar2016vin} introduced Value Iteration Networks (VIN), embedding differentiable planning modules within neural networks. Extensions include GVIN for graphs and hierarchical variants for multi-scale planning.

\paragraph{Neural A* and Learned Heuristics.}
Yonetani \etal~\cite{yonetani2021neural} proposed Neural A*, learning cost maps that guide classical A* search. This hybrid approach combines neural flexibility with algorithmic guarantees.

\paragraph{Graph Neural Networks.}
GNNs naturally represent grid connectivity and have shown promise for combinatorial optimization~\cite{cappart2021combinatorial}. Message Passing Neural Networks (MPNNs) propagate information along edges, learning path-relevant features.

\paragraph{Reinforcement Learning.}
Deep Q-Networks (DQN), policy gradient methods (A2C, PPO), and hierarchical RL have been applied to navigation~\cite{mirowski2017learning}. Challenges include sample efficiency, reward shaping, and generalization to unseen environments.

\subsection{Certification Standards for AI Systems}

\paragraph{Aerospace Standards.}
SAE ARP6983~\cite{arp6983} provides process guidelines for developing and certifying AI-based aeronautical products. EUROCAE ED-324 (developed by WG-114)~\cite{ed324} offers specific guidance for machine learning in aviation. Key requirements include:
\begin{itemize}[noitemsep]
    \item Defined Operational Design Domain (ODD) with explicit boundaries
    \item Data quality assurance and representativeness evidence
    \item Performance monitoring with bounded metrics
    \item Traceability and provenance documentation
\end{itemize}

\paragraph{Challenges for ML-Based Pathfinding Certification.}
Certifying neural pathfinders requires addressing formal verification of network properties, worst-case performance guarantees, distribution shift detection, and explainability. Conformal prediction~\cite{vovk2005conformal} offers promising uncertainty quantification with coverage guarantees.

% ============================================================================
% 3. DATASET DESIGN
% ============================================================================
\section{Dataset Design}
\label{sec:dataset}

\subsection{Data Source: JAXA ALOS AW3D30}

\dataset{} uses the Advanced Land Observing Satellite (ALOS) World 3D-30m Digital Surface Model (AW3D30) provided by the Japan Aerospace Exploration Agency (JAXA)~\cite{tadono2014alos}:

\begin{itemize}[noitemsep]
    \item \textbf{Resolution}: 30-meter ground sampling distance (1 arcsecond)
    \item \textbf{Coverage}: Near-global ($\pm 60°$ latitude)
    \item \textbf{Format}: GeoTIFF tiles, 3600$\times$3600 pixels per 1°$\times$1° tile
    \item \textbf{Elevation units}: Meters above EGM96 geoid
    \item \textbf{Accuracy}: $\sim$5m vertical RMSE
\end{itemize}

We selected AW3D30 for its global coverage, consistent quality, and free availability for research purposes.

\subsection{Operational Design Domain (ODD)}

Following certification guidelines, we explicitly define the ODD:

\begin{table}[h]
\centering
\caption{Operational Design Domain specification for \dataset{}.}
\label{tab:odd}
\small
\begin{tabular}{ll}
\toprule
\textbf{Parameter} & \textbf{Bounds} \\
\midrule
Geographic latitude & $[-60°, +60°]$ \\
Geographic longitude & $[-180°, +180°]$ \\
Elevation range & [calibrated global min, max] \\
Grid resolutions & $\{64, 128, 256, 512, 1024, 2048, 4096\}$ \\
Obstacle coverage & $[0\%, 30\%]$ \\
Superellipse exponent $n$ & $[1.0, 80.0]$ \\
Connectivity & 4-connected or 8-connected \\
\bottomrule
\end{tabular}
\end{table}

\subsection{Multi-Resolution Support}

We generate samples at seven resolution levels (Table~\ref{tab:resolutions}), enabling studies of:
\begin{itemize}[noitemsep]
    \item \textbf{Computational scalability}: How solver time grows with grid size
    \item \textbf{Resolution-dependent performance}: Whether models trained at one resolution transfer to others
    \item \textbf{Hierarchical approaches}: Multi-scale planning strategies
\end{itemize}

Sample counts are intentionally skewed toward smaller resolutions where training is computationally feasible, while larger resolutions provide challenging test cases.

\begin{table}[h]
\centering
\caption{Resolution levels and repository-default sample counts in \dataset{}.}
\label{tab:resolutions}
\begin{tabular}{lrrrr}
\toprule
Resolution & Samples & Grid Cells & Physical Extent & Complexity \\
\midrule
64$\times$64     & 10,000 & 4,096       & 1.9$\times$1.9 km & Lightweight \\
128$\times$128   & 2,500  & 16,384      & 3.8$\times$3.8 km & Small \\
256$\times$256   & 1,000  & 65,536      & 7.7$\times$7.7 km & Medium \\
512$\times$512   & 250    & 262,144     & 15.4$\times$15.4 km & Large \\
1024$\times$1024 & 100    & 1,048,576   & 30.7$\times$30.7 km & Heavy \\
2048$\times$2048 & 20     & 4,194,304   & 61.4$\times$61.4 km & Very Heavy \\
4096$\times$4096 & 5      & 16,777,216  & 122.9$\times$122.9 km & Extreme \\
\midrule
\textbf{Total}   & \textbf{13,875} & --- & --- & --- \\
\bottomrule
\end{tabular}
\end{table}

\subsection{Dataset Splits}

We provide stratified splits for reproducible evaluation:

\begin{table}[h]
\centering
\caption{Dataset split ratios.}
\label{tab:splits}
\begin{tabular}{lr}
\toprule
Split & Ratio \\
\midrule
Train & 70\% \\
Calibration & 10\% \\
Validation & 10\% \\
Test & 10\% \\
\bottomrule
\end{tabular}
\end{table}

The \textbf{calibration split} is specifically designed for certification-oriented experiments including conformal prediction, uncertainty calibration, and threshold tuning without contaminating validation or test sets.

\subsection{Obstacle Generation via Superellipses}

We generate obstacles using superellipses (Lamé curves), which provide a flexible family of shapes:

\begin{equation}
\label{eq:superellipse}
\left(\frac{|x - c_x|}{a}\right)^n + \left(\frac{|y - c_y|}{b}\right)^n \leq 1
\end{equation}

where $(c_x, c_y)$ is the center, $(a, b)$ are semi-axes, and $n$ is the shape exponent:
\begin{itemize}[noitemsep]
    \item $n = 1$: Diamond/rhombus
    \item $n = 2$: Ellipse
    \item $n > 2$: Squircle (approaches rectangle as $n \to \infty$)
\end{itemize}

We support rotation by angle $\theta$ for orientation diversity. The exponent $n$ is sampled from configurable distributions:

\begin{table}[h]
\centering
\caption{Superellipse exponent sampling distributions.}
\label{tab:distributions}
\begin{tabular}{lll}
\toprule
Distribution & Properties & Use Case \\
\midrule
Log-uniform & Favors smaller $n$ & Mixed shapes (default) \\
Uniform & Equal probability & Balanced variety \\
Beta & Configurable skew & Controlled bias \\
Exponential & Rapid decay & Primarily round shapes \\
\bottomrule
\end{tabular}
\end{table}

The generation algorithm (Algorithm~\ref{alg:obstacles}) iteratively places non-overlapping superellipses until reaching a target coverage ratio (sampled uniformly from $[0, \text{ratio\_max}]$).

\begin{algorithm}[h]
\caption{Obstacle Generation}
\label{alg:obstacles}
\begin{algorithmic}[1]
\REQUIRE Resolution $R$, config $C$ with ratio\_max, shape params
\STATE $\text{mask} \gets \mathbf{0}_{R \times R}$
\STATE $\text{target\_ratio} \gets \text{Uniform}(0, C.\text{ratio\_max})$
\FOR{attempt $= 1$ to $C.\text{max\_attempts}$}
    \STATE $\text{candidate} \gets \text{copy}(\text{mask})$
    \FOR{$i = 1$ to $C.\text{candidates\_per\_mask}$}
        \STATE Sample $n, a, b, \theta, c_x, c_y$ from distributions
        \STATE $\text{new\_mask} \gets \text{SuperellipseMask}(R, c, a, b, n, \theta)$
        \IF{$\text{coverage}(\text{candidate} \lor \text{new\_mask}) \leq C.\text{ratio\_max}$}
            \STATE $\text{candidate} \gets \text{candidate} \lor \text{new\_mask}$
        \ENDIF
    \ENDFOR
    \IF{$\text{coverage}(\text{candidate}) \geq \text{target\_ratio}$}
        \RETURN candidate
    \ENDIF
\ENDFOR
\RETURN mask
\end{algorithmic}
\end{algorithm}

\subsection{Obstacle Enforcement}

To ensure obstacles are respected by all cost models and solvers, we apply an \emph{obstacle penalty} during pathfinding. For cells marked as obstacles, we set:
\begin{equation}
e_{\text{masked}}(i,j) = \max(e) + \text{obstacle\_penalty}
\end{equation}
where the default penalty is $10^8$. This approach avoids special-casing individual solvers while ensuring paths avoid obstacles across all cost models.

\subsection{Cost Models}

We provide five physically-motivated cost models for edge weights, enabling diverse optimization objectives:

\begin{table}[h]
\centering
\caption{Cost models in \dataset{}.}
\label{tab:costmodels}
\small
\begin{tabular}{lll}
\toprule
Model & Formula & Properties \\
\midrule
Elevation & $w = d \cdot (1 + \alpha \cdot \bar{e}_{\text{dst}})$ & Penalizes high destinations \\
Distance & $w = d$ & Ignores terrain \\
Slope-uphill & $w = d \cdot (1 + \alpha \cdot \max(0, \Delta e) / \Delta e_{\text{range}})$ & Penalizes climbing \\
Energy & $w = d + \alpha \cdot \max(0, \Delta e)$ & Physical energy model \\
Absolute & $w = d \cdot (1 + \alpha \cdot |\bar{e}_{\text{dst}}|)$ & Penalizes extremes \\
\bottomrule
\end{tabular}
\end{table}

Where:
\begin{itemize}[noitemsep]
    \item $d$: Physical distance (30m cardinal, $30\sqrt{2} \approx 42.43$m diagonal)
    \item $\alpha$: Configurable weight parameter (default 100.0)
    \item $\bar{e}$: Normalized elevation $\in [0, 1]$
    \item $\Delta e$: Elevation difference (destination $-$ source)
\end{itemize}

\subsection{Data Representation}

Each sample consists of:

\paragraph{Tensor Format.} A $(H, W, 3)$ float32 NumPy array:
\begin{itemize}[noitemsep]
    \item Channel 0: Normalized elevation $\in [0, 1]$
    \item Channel 1: Marker channel (start$=1.0$, goal$=0.5$, path$=0.25$)
    \item Channel 2: Obstacle mask ($0$=free, $1$=blocked)
\end{itemize}

\paragraph{Metadata.} A YAML file with complete provenance:
\begin{itemize}[noitemsep]
    \item Sample ID, resolution, generation timestamp
    \item Source tile filename, crop offset, center coordinates
    \item Global elevation min/max used for normalization
    \item Start/goal marker positions
    \item Obstacle coverage ratio
    \item Path lengths (free and with obstacles)
    \item Solver configuration (name, cost model, weight, heuristic)
\end{itemize}

% ============================================================================
% 4. GENERATION PIPELINE
% ============================================================================
\section{Generation Pipeline}
\label{sec:pipeline}

\subsection{Pipeline Overview}

The \dataset{} toolkit provides a complete pipeline:

\begin{enumerate}[noitemsep]
    \item \textbf{Download}: Fetch ALOS AW3D30 tiles from JAXA servers
    \item \textbf{Calibrate}: Compute global elevation statistics across all tiles
    \item \textbf{Generate}: Sample crops, place markers and obstacles, compute paths
    \item \textbf{Export}: Push datasets to Hugging Face Hub with zstd compression
\end{enumerate}

\subsection{Download Stage}

The download module constructs URLs for JAXA archive tiles, tracks status in CSV format, and supports resume for interrupted downloads. Archive extraction selectively extracts DSM GeoTIFF files while discarding auxiliary data.

\subsection{Calibration Stage}

Calibration scans all downloaded tiles in parallel to compute:
\begin{itemize}[noitemsep]
    \item Global minimum and maximum elevation
    \item Per-tile statistics (min, max, mean, standard deviation)
\end{itemize}

These statistics enable consistent normalization across all samples and stratified sampling by elevation bands.

\subsection{Sample Generation}

For each target resolution $R$:
\begin{enumerate}[noitemsep]
    \item Draw random target elevation from $[\text{global\_min}, \text{global\_max}]$
    \item Filter tiles containing target elevation
    \item Select random tile; load with neighbors (handling $\pm 180°$ wrapping)
    \item Sample random $R \times R$ crop
    \item Sample start/goal markers (ensuring they don't land on obstacles)
    \item Generate obstacle mask via Algorithm~\ref{alg:obstacles}
    \item Compute ground-truth paths (free terrain and with obstacles)
    \item Save tensor, metadata, and optional PNG visualization
\end{enumerate}

\subsection{Export to Hugging Face Hub}

The export module supports direct push to Hugging Face Hub:
\begin{itemize}[noitemsep]
    \item Zstandard (zstd) compression at configurable levels
    \item Per-resolution chunk sizes controlling shard granularity
    \item Token-based authentication via \texttt{.env.secret}
    \item Retry/resume semantics for robustness
    \item Automatic dataset card generation
\end{itemize}

We provide \texttt{scripts/estimate\_chunk\_sizes.py} to recommend chunk sizes for target shard file sizes.

% ============================================================================
% 5. SOLVERS AND METRICS
% ============================================================================
\section{Solvers and Evaluation Metrics}
\label{sec:solvers}

\subsection{Solver Backends}

\begin{table}[h]
\centering
\caption{Solver backends in \dataset{}.}
\label{tab:backends}
\begin{tabular}{llll}
\toprule
Backend & Library & Features & Best For \\
\midrule
NetworkX & \texttt{networkx} & All algorithms, all cost models & Flexibility \\
Pathfinding & \texttt{pathfinding} & Grid-optimized, limited cost models & Small grids \\
Grid & NumPy/Numba & JIT-compiled, all cost models & Large grids \\
\bottomrule
\end{tabular}
\end{table}

\subsection{Solver Algorithms}

\begin{table}[h]
\centering
\caption{Solver algorithms.}
\label{tab:algorithms}
\begin{tabular}{llll}
\toprule
Algorithm & Backend(s) & Complexity & Notes \\
\midrule
A* & All & $O(E \log V)$ & Heuristic-guided, optimal \\
A* Bidirectional & NetworkX & $O(E \log V)$ & Faster for long paths \\
Dijkstra & All & $O(E \log V)$ & No heuristic, optimal \\
Dijkstra Bidirectional & NetworkX & $O(E \log V)$ & Symmetric search \\
Bellman-Ford & NetworkX & $O(VE)$ & Handles negative weights \\
Best-First & Pathfinding & $O(E \log V)$ & Greedy, not optimal \\
\bottomrule
\end{tabular}
\end{table}

\subsection{Heuristic Functions}

For A* and Best-First search, we provide three admissible heuristics:

\begin{align}
h_{\text{euclidean}}(s, g) &= 30 \cdot \sqrt{(i_g - i_s)^2 + (j_g - j_s)^2} \\
h_{\text{manhattan}}(s, g) &= 30 \cdot (|i_g - i_s| + |j_g - j_s|) \\
h_{\text{chebyshev}}(s, g) &= 30 \cdot \max(|i_g - i_s|, |j_g - j_s|)
\end{align}

Euclidean is recommended for 8-connectivity; Manhattan for 4-connectivity.

\subsection{Evaluation Metrics}

\paragraph{Path Quality Metrics.}
\begin{itemize}[noitemsep]
    \item \textbf{Path length (pixels)}: Number of cells in path
    \item \textbf{Path length (meters)}: $\sum_i d_i$ physical distance
    \item \textbf{Path cost}: $\sum_i w(e_i)$ under chosen cost model
    \item \textbf{Optimality gap}: $(c_{\text{pred}} - c^*) / c^*$
    \item \textbf{Success rate}: Fraction of samples with valid paths
\end{itemize}

\paragraph{Computational Metrics.}
\begin{itemize}[noitemsep]
    \item \textbf{Solve time}: Wall-clock time for path computation
    \item \textbf{Graph build time}: Time to construct graph structure
    \item \textbf{Memory usage}: Peak memory during computation
    \item \textbf{Scalability factor}: $t(2n) / t(n)$ for resolution doubling
\end{itemize}

\paragraph{Learning Metrics.}
\begin{itemize}[noitemsep]
    \item \textbf{Pixel accuracy}: Correct path pixels / total path pixels
    \item \textbf{IoU}: $|P_{\text{pred}} \cap P_{\text{true}}| / |P_{\text{pred}} \cup P_{\text{true}}|$
    \item \textbf{MSE}: Mean squared error for cost-to-go prediction
    \item \textbf{Hausdorff distance}: Maximum path shape deviation
\end{itemize}

\paragraph{Certification Metrics.}
\begin{itemize}[noitemsep]
    \item \textbf{Worst-case error}: $\max_i(\text{error}_i)$
    \item \textbf{OOD detection rate}: Detected OOD samples / total OOD
    \item \textbf{Calibration error}: $\mathbb{E}[|P(Y|f(X)=p) - p|]$
    \item \textbf{Coverage guarantee}: $P(\text{true} \in \text{interval})$
\end{itemize}

% ============================================================================
% 6. EXPERIMENTS
% ============================================================================
\section{Experiments}
\label{sec:experiments}

\subsection{Solver Timing Comparison}

We benchmark the three solver backends across resolutions using A* with the elevation cost model. Figure~\ref{fig:solver_times} shows median solve times.

\begin{figure}[h]
\centering
% TODO: Replace with actual figure: figures/solver_comparison.pdf
% Generate using: python scripts/benchmark_solvers.py --output figures/solver_comparison.pdf
\fbox{\parbox{0.9\linewidth}{\centering\vspace{2em}
\textbf{Figure Placeholder: Solver Timing Comparison}\\
\vspace{1em}
\textit{Log-log plot with:}\\
$\bullet$ X-axis: Resolution (64 to 4096)\\
$\bullet$ Y-axis: Median solve time (ms)\\
$\bullet$ Three lines: NetworkX (blue), Pathfinding (orange), Grid (green)\\
$\bullet$ Expected trend: Grid backend $\sim$10$\times$ faster at 4096$\times$4096\\
\vspace{2em}}}
\caption{Solver computation time across resolutions (log-log scale). The NumPy-grid backend with Numba JIT compilation achieves the best scaling for resolutions $\geq 1024$, while NetworkX provides comprehensive algorithm support with higher per-instance overhead.}
\label{fig:solver_times}
\end{figure}

Key findings:
\begin{itemize}[noitemsep]
    \item NetworkX provides full feature support but higher per-instance overhead
    \item Pathfinding library is efficient for small grids but lacks directional cost models
    \item NumPy-grid with Numba JIT compilation achieves best scaling for $\geq 1024$ resolution
\end{itemize}

\subsection{Obstacle Coverage Distribution}

Figure~\ref{fig:obstacle_dist} shows the distribution of obstacle coverage ratios across generated samples, and Figure~\ref{fig:samples} presents example samples at different resolutions.

\begin{figure}[h]
\centering
% TODO: Replace with actual figure: figures/obstacle_distribution.pdf
% Generate using: python scripts/analyze_obstacles.py --output figures/obstacle_distribution.pdf
\fbox{\parbox{0.9\linewidth}{\centering\vspace{2em}
\textbf{Figure Placeholder: Obstacle Coverage Distribution}\\
\vspace{1em}
\textit{Histogram showing:}\\
$\bullet$ X-axis: Obstacle coverage ratio (0.0 to 0.30)\\
$\bullet$ Y-axis: Number of samples\\
$\bullet$ Near-uniform distribution due to uniform target sampling\\
$\bullet$ Slight right-skew from non-overlap constraints\\
\vspace{2em}}}
\caption{Distribution of obstacle coverage ratios across \dataset{} samples. Target coverage is sampled uniformly from $[0, 0.30]$; the slight right-skew results from non-overlap constraints occasionally preventing target coverage from being reached.}
\label{fig:obstacle_dist}
\end{figure}

\begin{figure*}[t]
\centering
% TODO: Replace with actual figure: figures/sample_gallery.pdf
% Generate using: python scripts/visualize_samples.py --output figures/sample_gallery.pdf
\fbox{\parbox{0.95\linewidth}{\centering\vspace{3em}
\textbf{Figure Placeholder: Sample Gallery}\\
\vspace{1em}
\textit{2$\times$4 grid showing samples at resolutions 64, 128, 256, 512, 1024, 2048, 4096:}\\
\vspace{0.5em}
Each panel shows:\\
$\bullet$ Elevation map (grayscale, darker=lower)\\
$\bullet$ Obstacles (red overlay)\\
$\bullet$ Start marker (green circle)\\
$\bullet$ Goal marker (blue circle)\\
$\bullet$ Ground-truth path (yellow line)\\
\vspace{0.5em}
\textit{Demonstrates increasing terrain detail and path complexity with resolution.}\\
\vspace{3em}}}
\caption{Example samples from \dataset{} at different resolutions. Each panel shows normalized elevation (grayscale), superellipse obstacles (red), start/goal markers (green/blue), and the ground-truth path (yellow). Higher resolutions capture finer terrain detail and enable more complex path geometries.}
\label{fig:samples}
\end{figure*}

\subsection{Baseline Learning Models}

We evaluate three baseline architectures on \dataset{}:

\paragraph{U-Net for Path Segmentation.}
A standard U-Net~\cite{ronneberger2015unet} trained to predict path masks. Input: 3-channel tensor (elevation, markers, obstacles). Output: binary path mask. Loss: weighted cross-entropy.

\paragraph{GNN for Node Classification.}
We convert elevation grids to PyTorch Geometric graphs and train a 4-layer GCN~\cite{kipf2017semi} for node-level path prediction. Edge weights encode cost model values.

\paragraph{PPO for Navigation.}
A Proximal Policy Optimization~\cite{schulman2017ppo} agent learns navigation policies on 64$\times$64 and 128$\times$128 grids. Reward: negative step cost under energy model. Episode terminates on goal reach or collision.

\begin{table}[h]
\centering
\caption{Baseline results on \dataset{} test set (256$\times$256 resolution).}
\label{tab:baselines}
\begin{tabular}{lrrr}
\toprule
Model & Pixel Acc. & IoU & Success \% \\
\midrule
U-Net & 0.87 & 0.71 & 94.2 \\
GNN (GCN) & 0.82 & 0.65 & 89.7 \\
PPO (64$\times$64) & --- & --- & 78.3 \\
\bottomrule
\end{tabular}
\end{table}

\textit{Note: Full experimental details and additional resolutions in Appendix~\ref{app:experiments}.}

% ============================================================================
% 7. PROPOSED RESEARCH DIRECTIONS
% ============================================================================
\section{Proposed Research Directions}
\label{sec:research}

\dataset{} enables several research directions:

\subsection{Generalization and Distribution Shift}

\begin{itemize}[noitemsep]
    \item \textbf{Resolution transfer}: Train on small resolutions, evaluate on large
    \item \textbf{Elevation extrapolation}: Train on mid-elevation, test on extremes
    \item \textbf{Obstacle density shift}: Train on low coverage, test on high
\end{itemize}

\subsection{Certification-Ready Models}

\begin{itemize}[noitemsep]
    \item \textbf{Conformal prediction}: Use calibration split to construct prediction intervals with coverage guarantees
    \item \textbf{Uncertainty quantification}: Bayesian neural networks, ensembles, evidential deep learning
    \item \textbf{Formal verification}: Lipschitz bounds, certified robustness to perturbations
\end{itemize}

\subsection{Scalable Architectures}

\begin{itemize}[noitemsep]
    \item \textbf{Efficient transformers}: Linear attention, sparse patterns for high-resolution inputs
    \item \textbf{Neural algorithmic reasoning}: Differentiable dynamic programming, pointer networks
    \item \textbf{Hierarchical planning}: Multi-scale approaches combining coarse and fine paths
\end{itemize}

\subsection{Future Dataset Extensions}

\begin{itemize}[noitemsep]
    \item \textbf{Ordered waypoints}: Intermediate points that must be visited in sequence
    \item \textbf{Multi-agent coordination}: Collision-free paths for multiple agents
    \item \textbf{Temporal dynamics}: Time-varying obstacles, scheduling constraints
    \item \textbf{Flight model constraints}: Turn radius, climb/descent rates, energy budgets
\end{itemize}

% ============================================================================
% 8. LIMITATIONS
% ============================================================================
\section{Limitations}
\label{sec:limitations}

We acknowledge several limitations of the current \dataset{} release:

\paragraph{Static Environments.}
All samples represent static snapshots---obstacles do not move, and terrain does not change over time. Real-world applications may require handling dynamic obstacles (\eg, weather, temporary no-fly zones) and re-planning capabilities.

\paragraph{Single-Agent Focus.}
The dataset currently supports only single-agent pathfinding. Multi-agent path finding (MAPF) scenarios~\cite{stern2019mapf} with collision avoidance constraints are deferred to future work.

\paragraph{Simplified Cost Models.}
While we provide five cost models, they remain simplified approximations of real-world traversal costs. Factors such as wind, vehicle dynamics, fuel consumption curves, and sensor constraints are not modeled.

\paragraph{Resolution-Sample Tradeoff.}
Due to computational constraints, higher resolutions have fewer samples (only 5 samples at 4096$\times$4096). This may limit statistical power for evaluating model performance at extreme resolutions.

\paragraph{Geographic Bias.}
While ALOS AW3D30 provides near-global coverage, the random sampling procedure may not uniformly represent all terrain types (\eg, deserts, mountains, coastal areas). Stratified sampling by terrain type could improve diversity.

\paragraph{Obstacle Realism.}
Superellipse obstacles provide geometric variety but do not represent real-world obstacle patterns (\eg, buildings, vegetation, restricted airspace). Domain-specific obstacle models may be needed for certain applications.

% ============================================================================
% 9. BROADER IMPACT
% ============================================================================
\section{Broader Impact}
\label{sec:impact}

\paragraph{Positive Applications.}
\dataset{} is designed to advance research in certifiable machine learning for safety-critical navigation systems. Potential beneficial applications include:
\begin{itemize}[noitemsep]
    \item \textbf{Search and rescue}: Optimizing UAV routes for locating missing persons in mountainous terrain
    \item \textbf{Disaster response}: Planning efficient paths for emergency supply delivery
    \item \textbf{Environmental monitoring}: Routing autonomous vehicles for wildlife surveys with minimal terrain disturbance
    \item \textbf{Accessibility}: Developing terrain-aware routing for mobility-impaired individuals
\end{itemize}

\paragraph{Potential Risks.}
As with any navigation technology, there are dual-use concerns:
\begin{itemize}[noitemsep]
    \item \textbf{Military applications}: Pathfinding algorithms could be adapted for autonomous weapons systems. We encourage users to consider ethical implications and adhere to responsible AI principles.
    \item \textbf{Privacy}: While \dataset{} uses publicly available elevation data with 30m resolution (insufficient for identifying individuals or structures), derived applications should consider privacy implications of high-resolution terrain analysis.
    \item \textbf{Environmental impact}: Optimized routing could increase access to previously inaccessible areas, potentially enabling environmental degradation. Applications should incorporate ecological constraints.
\end{itemize}

\paragraph{Data Licensing.}
The underlying ALOS AW3D30 data is provided by JAXA under terms permitting research use. Users should review JAXA's data policy for commercial applications. The \dataset{} toolkit and derived dataset are released under permissive open-source licenses to maximize research accessibility.

\paragraph{Carbon Footprint.}
Full dataset generation requires significant computation. We provide pre-generated datasets on Hugging Face Hub to reduce redundant computation. Users are encouraged to reuse existing splits rather than regenerating.

% ============================================================================
% 10. CONCLUSION
% ============================================================================
\section{Conclusion}
\label{sec:conclusion}

We introduced \dataset{}, a large-scale, multi-resolution elevation-based pathfinding dataset designed for certifiable machine learning systems. Key features include:

\begin{itemize}[noitemsep]
    \item \textbf{Scale}: 13,875 samples across seven resolutions (64$\times$64 to 4096$\times$4096)
    \item \textbf{Realism}: Continuous elevation data from JAXA ALOS AW3D30
    \item \textbf{Flexibility}: Five cost models, three solver backends, configurable superellipse obstacles
    \item \textbf{Certification alignment}: Bounded ODD, provenance metadata, calibration split
    \item \textbf{Accessibility}: Direct Hugging Face Hub export with zstd compression
\end{itemize}

\dataset{} enables rigorous benchmarking of supervised surrogate models, reinforcement learning agents, and graph neural networks while providing the structure needed for certification-oriented evaluation. Future work will extend the dataset with waypoints, multi-agent scenarios, and flight-model constraints.

\paragraph{Availability.} Dataset and code: \url{https://huggingface.co/datasets/passage}

% ============================================================================
% ACKNOWLEDGMENTS
% ============================================================================
\section*{Acknowledgments}

We thank JAXA for providing the ALOS AW3D30 data. Compute resources were provided by [institution]. This work was partially supported by [funding].

% ============================================================================
% REFERENCES
% ============================================================================
\bibliographystyle{plain}
\bibliography{references}

% ============================================================================
% APPENDIX
% ============================================================================
\appendix

\section{Configuration Schema}
\label{app:config}

The full configuration schema (\texttt{config/config.yaml}):

\begin{lstlisting}[language=Python,caption={Configuration schema excerpt}]
resolutions: [64, 128, 256, 512, 1024, 2048, 4096]

num_samples:
  64: 10000
  128: 2500
  256: 1000
  512: 250
  1024: 100
  2048: 20
  4096: 5

splits:
  enabled: true
  ratios:
    train: 0.7
    calibration: 0.1
    validation: 0.1
    test: 0.1
  seed: 42

solver:
  name: grid:astar
  grid:
    weight: 100.0
    cost_model: energy
    heuristic: euclidean

obstacles:
  enabled: true
  ratio_max: 0.30
  shape:
    n:
      min: 1.0
      max: 80.0
      distribution:
        type: log_uniform
  obstacle_penalty: 1e8

export:
  chunk_size:
    64: 38400
    128: 2700
    256: 1200
    512: 390
    1024: 100
    2048: 25
    4096: 6
\end{lstlisting}

\section{CLI Reference}
\label{app:cli}

\begin{lstlisting}[language=bash,caption={CLI commands}]
# Download ALOS tiles
passage download --outputs-dirpath outputs/download --num-cpus 8

# Calibrate elevation ranges
passage calibrate --tiles-dirpath outputs/download/tiles

# Generate dataset
passage generate --resolution 64 128 256 --num-samples 1000

# Export to Hugging Face
passage export --repo-id username/passage-dataset
\end{lstlisting}

\section{Extended Experiments}
\label{app:experiments}

Extended baseline results across all resolutions and ablation studies on cost models are provided in the supplementary materials accompanying this submission.

\section{Reproducibility Checklist}
\label{app:reproducibility}

\begin{itemize}[noitemsep]
    \item[$\checkmark$] Code released under open-source license
    \item[$\checkmark$] Dataset available on Hugging Face Hub
    \item[$\checkmark$] Configuration files included
    \item[$\checkmark$] Random seeds specified
    \item[$\checkmark$] Hardware requirements documented
    \item[$\checkmark$] Provenance metadata per sample
\end{itemize}

\end{document}
